% !TEX program = pdflatex
\documentclass[11pt,a4paper]{article}

\usepackage[T1]{fontenc}
\usepackage[utf8]{inputenc}
\usepackage{lmodern}
\usepackage{microtype}
\usepackage[margin=1in]{geometry}
\usepackage{amsmath,amssymb,amsthm}
\usepackage{mathtools}
\usepackage{booktabs}
\usepackage{array}
\usepackage{enumitem}
\usepackage{xcolor}
\usepackage{listings}
\usepackage{tabularx}
\usepackage[hidelinks,colorlinks=true,linkcolor=blue!50!black,citecolor=blue!50!black,urlcolor=blue!50!black]{hyperref}
\usepackage{cleveref}

% Allow long underscore-laden identifiers in \texttt{} to break across lines.
% Without this, names like crypto_core_ristretto255_is_valid_point overflow.
\makeatletter
\def\BreakableUnderscore{\leavevmode\nobreak\hskip\z@skip
  \discretionary{\textunderscore\hskip\z@skip}{\hskip\z@skip}{\textunderscore}%
  \nobreak\hskip\z@skip\relax}
\makeatother
\newcommand{\codename}[1]{\texttt{\detokenize{#1}\relax}\@ifnextchar.{\relax}{\relax}}
% Simpler: just enable hyphenation inside \texttt for our purposes.
\renewcommand{\sfdefault}{lmss}
\sloppy

% ---- macros ----------------------------------------------------------------
\newcommand{\G}{\mathbb{G}}
\newcommand{\Zq}{\mathbb{Z}_q}
\newcommand{\sk}{\mathit{sk}}
\newcommand{\pk}{\mathit{pk}}
\newcommand{\Hzero}{H_0}
\newcommand{\Hone}{H_1}
\newcommand{\Htwo}{H_2}
\newcommand{\dst}{\mathsf{DST}}
\newcommand{\sample}{\xleftarrow{\$}}
% Algorithm names usable in both text and math mode.
\newcommand{\Sign}{\textsf{Sign}}
\newcommand{\Verify}{\textsf{Verify}}
\newcommand{\Trace}{\textsf{Trace}}
\newcommand{\KeyGen}{\textsf{KeyGen}}

\theoremstyle{definition}
\newtheorem{assumption}{Assumption}
\newtheorem{remark}{Remark}

% ---- listings --------------------------------------------------------------
\definecolor{codebg}{rgb}{0.97,0.97,0.97}
\lstdefinestyle{shell}{
  basicstyle=\ttfamily\small,
  backgroundcolor=\color{codebg},
  frame=single,
  framesep=4pt,
  rulecolor=\color{black!15},
  breaklines=true,
  breakatwhitespace=true,
  postbreak=\mbox{\textcolor{gray}{$\hookrightarrow$}\space},
  showstringspaces=false,
}
\lstdefinestyle{py}{
  language=Python,
  basicstyle=\ttfamily\small,
  backgroundcolor=\color{codebg},
  frame=single,
  framesep=4pt,
  rulecolor=\color{black!15},
  keywordstyle=\color{blue!60!black}\bfseries,
  stringstyle=\color{green!30!black},
  commentstyle=\color{gray}\itshape,
  showstringspaces=false,
  breaklines=true,
}

% ---- title -----------------------------------------------------------------
\title{\textbf{Verifiable Anonymous Voting on Top of One-Time Traceable Ring Signatures:\\ Two Architectures and a Comparison}\\[0.5em]
       \large Research Artifact, v0.4}
\author{Daniele Lin \and Niccol\`o Pagano}
\date{May 2026}

\begin{document}
\maketitle

\begin{abstract}
We present a verifiable anonymous voting artifact built on Scafuro and
Zhang's one-time traceable ring signature scheme
(OTRS)~\cite{ScafuroZhang2021} and instantiated over Ristretto255 with
RFC~9380 hash-to-curve. On top of that single cryptographic primitive
we implement and compare \emph{two} system-layer architectures.
\textbf{Architecture~A} is a hash-chained append-only bulletin board
whose entries are co-signed by a $t$-of-$N$ Ed25519 publisher cohort and
checkpointed by a $k$-of-$M$ witness federation that surfaces cohort
equivocation publicly --- a federated, Certificate-Transparency-style
system. \textbf{Architecture~B} is a Bitcoin-shaped proof-of-work chain
with permissionless mining, an eVote ledger, typed voting transactions,
and the same OTRS-based ballot semantics. Both architectures support
the same election lifecycle (Setup $\to$ Registration$^*$ $\to$ Ring
$\to$ Ballot$^*$ $\to$ Closed $\to$ Tally?) and reuse the same record
schema. The artifact ships 127 tests --- including 16 covering the
PoW path end-to-end --- two CLIs, and a documented threat model.
We compare the architectures on trust model, decentralisation, cost,
censorship resistance, and cross-poll linkability, and argue that the
two answer different questions: A trades permissionless validator
membership for cheap, auditable accountability against a small named
authority set; B trades that for open mining at the cost of bringing
back the security-budget asymmetry that makes blockchain voting
fragile~\cite{ParkSpecter2021}. Residual gaps --- coercion resistance,
true threshold-Ed25519 (FROST), formal verification, post-quantum
variants --- are flagged for future work.
\end{abstract}

\section{Introduction}

Verifiable anonymous voting has become a recurring research theme because
nation-scale deployments still rely on trust-the-server
architectures whose properties are not publicly checkable~\cite{Springall2014}.
The cryptographic literature has long offered alternatives---Helios~\cite{Helios2008},
Civitas~\cite{Civitas2008}, end-to-end verifiable schemes such as
Scantegrity~\cite{Scantegrity2008} and, more recently,
ElectionGuard~\cite{ElectionGuard2021}---but each makes architectural
commitments that constrain the kinds of elections they fit.

Ring signatures, introduced by Rivest, Shamir and Tauman~\cite{RST01}, offer
a particularly appealing primitive for ballot anonymity: a signer proves
membership in a ring of public keys without revealing which member they are.
The challenge for voting is that a \emph{plain} ring signature lets a voter
sign repeatedly. Two parallel literatures address this: \emph{linkable} ring
signatures~\cite{LWW04,CLSAG2020} which expose a tag identifying re-use of
the same key, and \emph{traceable} ring signatures~\cite{FS07,ScafuroZhang2021}
which expose the actual public key on misuse.

This artifact implements Scafuro and Zhang's 2021
scheme~\cite{ScafuroZhang2021}. We chose it because (i) its trace algorithm
is deterministic and one-time-tag-based, not interactive; (ii) it requires no
trusted setup; (iii) the algebra is a clean $\Sigma$-OR proof over a single
prime-order group, which makes it a realistic target for both implementation
and future formal verification.

\subsection{Contributions}

\begin{enumerate}[leftmargin=*]
\item \textbf{OTRS reference implementation} (\texttt{otrs/}, $\approx$1\,100
      LOC Python) of \KeyGen, \Sign, \Verify, \Trace over Ristretto255, with
      all randomness drawn from \texttt{os.urandom} and all hashing routed
      through a domain-separated RFC~9380 hash-to-curve.
\item \textbf{Architecture A --- federated bulletin board}
      (\texttt{voting/}, $\approx$1\,800 LOC): a publisher-signed,
      hash-chained, append-only log; typed records for the election
      lifecycle ($\textsf{Setup} \to \textsf{Registration}^* \to
      \textsf{Ring} \to \textsf{Ballot}^* \to \textsf{Closed} \to
      \textsf{Tally}$); a $t$-of-$N$ publisher cohort with asynchronous
      threshold co-signing; a $k$-of-$M$ witness federation that
      checkpoints log heads and makes cohort equivocation publicly
      detectable; manager, voter, witness, and auditor principal APIs;
      and a CLI for the full flow.
\item \textbf{Architecture B --- proof-of-work chain}
      (\texttt{chain/}, $\approx$1\,400 LOC): a Bitcoin-shaped SHA-256
      PoW chain with heaviest-chain selection and median-time difficulty
      adjustment; an account ledger with nonces, fees, and a coinbase
      reward; eight typed transactions (coinbase, transfer, plus
      ($\textsf{setup\_poll}, \textsf{register\_voter},
      \textsf{publish\_ring}, \textsf{ballot}, \textsf{close\_poll},
      \textsf{tally})$ that wrap the same record schema as
      Architecture~A); a permissionless miner; a chain-aware auditor;
      and a CLI parallel to A's.
\item \textbf{Comparison} of the two architectures on trust model,
      decentralisation, cost, censorship resistance, and cross-poll
      linkability (\cref{sec:compare}).
\item \textbf{Threat model} (\texttt{paper/threat\_model.md}) covering
      both architectures.
\item \textbf{Test suite} (\texttt{tests/}, 127 tests): algebraic
      invariants of the group wrapper, hash-to-curve determinism and
      domain separation, OTRS sign/verify correctness, negative cases,
      property-based tests under Hypothesis, bulletin-board integrity,
      threshold cohort and witness behaviour, end-to-end Architecture~A
      election flow, adversarial cases (double-sign, forged attestation,
      tampered ballot, wrong publisher key, extra records, ballot
      outside the ring, illegal choice), and end-to-end Architecture~B
      blockchain flow (PoW, ledger arithmetic, nonce replay, tx
      signature tampering, on-chain double-sign detection,
      save/load round-trip).
\item \textbf{Microbenchmark} (\texttt{bench/}) for the cryptographic
      primitive (\cref{sec:eval}).
\item \textbf{Critique and revival} of the prior implementation in
      \texttt{legacy/} --- a 160-bit DSA-style subgroup with
      \texttt{random.randrange} randomness, no tests, and a half-built
      PoW ledger. Architecture~B is in part the legacy ambition
      modernised: same PoW backbone, hardened crypto, and the missing
      pieces (eVote ledger, typed voting transactions, audit pipeline)
      filled in.
\end{enumerate}

\subsection{Non-goals}

We do not claim a side-channel-hardened implementation: Python's \texttt{int}
arithmetic is not constant-time, and our wrapper around libsodium has a
small number of branches on potentially derived scalars (always with
negligible collision probability). We also do not deliver a formal proof in
EasyCrypt or ProVerif; both are flagged as next steps (\cref{sec:future}).

\section{Preliminaries}

We work in a prime-order cyclic group $\G$ of prime order $q$ with generator
$g$, in which the Decisional Diffie-Hellman (DDH) problem is assumed hard.
We instantiate $\G$ as Ristretto255~\cite{Ristretto2020}, a prime-order
group of order
\[ q = 2^{252} + 27742317777372353535851937790883648493 \]
built on top of Curve25519. Ristretto eliminates the small-subgroup and
cofactor pitfalls of raw Edwards25519, ruling out a class of implementation
bugs by construction.

We require three hash functions modelled as random oracles:
\[
  \Hzero : \{0,1\}^* \to \G, \quad
  \Hone  : \{0,1\}^* \to \G, \quad
  \Htwo  : \{0,1\}^* \to \Zq.
\]
Both $\Hzero, \Hone$ are instantiated as the RFC~9380~\cite{RFC9380}
random-oracle hash-to-curve for Ristretto255 with the suite identifier
\texttt{ristretto255\_XMD:SHA-512\_R255MAP\_RO\_}. $\Htwo$ uses the same XMD
expander followed by \texttt{crypto\_core\_ristretto255\_scalar\_reduce}
(uniform on $\Zq$ within statistical distance $2^{-126}$).

Per RFC~9380 best practice we tag each oracle with a distinct
domain-separation string fixed in \texttt{otrs/otrs.py}:
\begin{lstlisting}[style=py]
DST_H0 = b"OTRS-v1-RISTRETTO255-XMD:SHA-512-H0"
DST_H1 = b"OTRS-v1-RISTRETTO255-XMD:SHA-512-H1"
DST_H2 = b"OTRS-v1-RISTRETTO255-XMD:SHA-512-H2"
\end{lstlisting}
Any change to the algebraic structure (curve, hash, transcript format)
\emph{must} bump the version prefix.

\begin{assumption}[DDH in $\G$]
For uniformly random $a, b, c \in \Zq$, the distributions
$(g^a, g^b, g^{ab})$ and $(g^a, g^b, g^c)$ are computationally indistinguishable.
\end{assumption}

\section{The Scafuro--Zhang construction}
\label{sec:construction}

Fix an \emph{issue} $I$---in the voting application, the election
identifier---and a ring $R = (\pk_1, \ldots, \pk_n)$ of $n$ public keys. We
use $1$-indexed positions so we can take modular inverses; $i \in \{1, \ldots, n\}$
is invertible in $\Zq$ since $n \ll q$.

\paragraph{\KeyGen.}
$x \sample \Zq^*$, $\pk = g^x$.

\paragraph{$\Sign(x_i, I, m, R, i)$.}
\begin{enumerate}[leftmargin=*]
\item Compute $h \gets \Hzero(I \,\|\, R)$ and $A_0 \gets \Hone(I \,\|\, R \,\|\, m)$.
\item Compute the trace tag $\sigma_i \gets h^{x_i}$ (this is what makes
      the scheme one-time-traceable: $\sigma_i$ depends only on $(I, R, x_i)$,
      not on $m$).
\item Set $A_1 \gets (\sigma_i / A_0)^{1/i}$. For $j \ne i$, define
      $\sigma_j \gets A_0 \cdot A_1^{\,j}$; observe $\sigma_i$ also satisfies
      the recurrence by construction.
\item For $j \ne i$ sample $c_j, z_j \sample \Zq$ and set
      $a_j \gets g^{z_j} \cdot \pk_j^{c_j}$,
      $b_j \gets h^{z_j} \cdot \sigma_j^{c_j}$.
\item For $j = i$ sample $w \sample \Zq$ and set $a_i \gets g^w$, $b_i \gets h^w$.
\item Compute
      \[c^* \gets \Htwo(I \,\|\, R \,\|\, A_0 \,\|\, A_1 \,\|\, a_1 \,\|\, \cdots \,\|\, a_n \,\|\, b_1 \,\|\, \cdots \,\|\, b_n).\]
\item Set $c_i \gets c^* - \sum_{j \ne i} c_j$ and $z_i \gets w - c_i x_i$.
\item Output $\sigma \gets (A_1, c_1, \ldots, c_n, z_1, \ldots, z_n)$.
\end{enumerate}

\paragraph{$\Verify(\sigma, I, m, R)$.}
Recompute $h$, $A_0$, $\sigma_j = A_0 \cdot A_1^j$,
$a_j = g^{z_j} \cdot \pk_j^{c_j}$,
$b_j = h^{z_j} \cdot \sigma_j^{c_j}$ for all $j$, and the challenge $c^*$ as
above. Accept iff $\sum_j c_j \equiv c^* \pmod{q}$.

\paragraph{$\Trace(R, I, (m_1, \sigma^{(1)}), (m_2, \sigma^{(2)}))$.}
Recompute $h$, then $A_0^{(k)} = \Hone(I \,\|\, R \,\|\, m_k)$ for $k \in \{1,2\}$,
and $\sigma_j^{(k)} = A_0^{(k)} \cdot (A_1^{(k)})^j$ for $j = 1, \ldots, n$.
For the \emph{real} signer at position $i^*$ we have
$\sigma_{i^*}^{(k)} = h^{x_{i^*}}$ for both $k$ (the trace tag is
$m_k$-independent), so column $i^*$ matches. For $j \ne i^*$ the tag depends
on $m_k$, so columns generically differ. Outcome:

\begin{itemize}[leftmargin=*]
\item \textbf{exactly one} column matches $\Rightarrow$ that column is the double-signer;
\item \textbf{all} columns match $\Rightarrow$ the two signatures are
      identical (replay or signature on the same message twice---we label
      this ``linked'');
\item \textbf{no} columns match $\Rightarrow$ independent signers.
\end{itemize}

This is implemented in \texttt{otrs.otrs.trace}.

\section{Implementation}
\label{sec:impl}

\subsection{Layout}
\begin{lstlisting}[style=shell]
otrs/
  group.py        # Ristretto255 wrapper: Scalar, Point, base_mul, ORDER
  hash.py         # RFC 9380 XMD:SHA-512 expand + hash-to-curve, hash-to-scalar
  otrs.py         # keygen / sign / verify / trace
  serialize.py    # canonical ring + signature encodings
  _libsodium.py   # direct cffi binding to libsodium.so.23
tests/            # unit, property, serialization tests
bench/            # microbenchmarks
paper/            # this document
legacy/           # the prior implementation, preserved for reference
\end{lstlisting}

\subsection{Group wrapper}
\texttt{otrs.group} exposes Ristretto255 through \texttt{Scalar} and
\texttt{Point} frozen dataclasses with byte-level canonical encodings
(32~bytes each). All operations delegate to libsodium via a direct cffi
ABI binding (\texttt{otrs/\_libsodium.py}); we did not rely on PyNaCl
because the Ubuntu-shipped build does not export Ristretto255 bindings.
The \texttt{Point} constructor calls
\mbox{\texttt{crypto\_core\_}}\allowbreak\mbox{\texttt{ristretto255\_}}\allowbreak\mbox{\texttt{is\_valid\_point}}
so caller-supplied encodings cannot bypass group-membership checks.

\subsection{Hash-to-curve}
We implement \texttt{expand\_message\_xmd} (RFC~9380~\S5.3.1) over SHA-512
in pure Python ($\approx$25 lines) and compose it with
\texttt{crypto\_core\_ristretto255\_from\_hash} to obtain the random-oracle
hash-to-curve. The same expander, followed by
\texttt{crypto\_core\_ristretto255\_scalar\_reduce}, gives a
uniform-by-statistical-distance hash-to-scalar. The DST-oversize path
(\texttt{H2C-OVERSIZE-DST-} prefix, RFC~9380~\S5.3.3) is implemented and
tested.

\subsection{Randomness}
All secret randomness flows from \texttt{os.urandom} via
\texttt{Scalar.random()}, which pulls 64~bytes and reduces $\bmod{}\, q$.
This is the only RNG path in the artifact---there is no rejection-sampled
custom PRG, no \texttt{random.randrange}, no seedable fallback. The legacy
implementation, by contrast, sampled keys with \texttt{random.randrange}
(predictable from a Mersenne-Twister state) and implemented a hash-based
PRG that is incompatible with the random-oracle assumption.

\subsection{Canonical encoding}
The transcript hashed into $\Htwo$ depends bit-for-bit on the encoding of
the ring and the auxiliary points $(a_j, b_j)$. We commit to the encoding
in \texttt{otrs/serialize.py}: a 4-byte length prefix, then concatenated
32-byte Ristretto encodings. Order is significant, so permuting the ring
changes the signature; this is required for the OR-proof structure.

\section{System architecture}
\label{sec:system}

The cryptographic primitive in \cref{sec:construction} gives us
\emph{ballot anonymity} and \emph{double-vote detection}. Turning that
into a usable election system requires three additional pieces: a
tamper-evident public record of everything that happened, a typed schema
for what can validly appear on that record, and well-defined roles for
the principals who interact with it. \texttt{voting/} delivers these.

\subsection{Bulletin board}
\label{sec:bb}

The bulletin board is an \emph{append-only, hash-chained, cohort-signed}
log. Each entry has the shape

\[ \mathsf{Entry} = (i, h_{\mathit{prev}}, t, \mathit{payload},
   \{(k_1, \sigma_1), \ldots, (k_s, \sigma_s)\}) \]

where $i$ is a contiguous index starting at $0$, $h_{\mathit{prev}}$ is
the SHA-256 hash of entry $i-1$ (or all-zero for the genesis entry), $t$
is a Unix timestamp, $\mathit{payload}$ is opaque bytes interpreted by
the record schema below, and the final set is a collection of
\emph{cohort-member signatures}: each $\sigma_j$ is an Ed25519 signature
by cohort member $k_j$ over the canonical content preimage. An entry is
valid only when at least $t$ distinct cohort members have signed it,
where $t$ is the threshold pinned in the genesis :class:`ElectionSetup`.
The entry hash is computed over the content alone, so collecting
additional signatures does not perturb the chain.

This is the Certificate-Transparency / Sigsum / C2SP transparency-log
model rather than a proof-of-work blockchain. Elections do not face
Sybil-mining attacks; they face misbehaving administrators. A
threshold-signed log is simpler, audited, and fits the actual threat
model. The cohort layer gives \emph{liveness against $N - t$
unavailable publishers} and \emph{soundness against $t - 1$ corrupted
publishers}. Equivocation \emph{within} the cohort
($\geq t$ corrupted) is addressed by witnesses (\cref{sec:witness}).

\subsection{Record schema and state machine}
\label{sec:records}

Each payload is a canonical JSON object discriminated by a \texttt{kind}
field. Six record kinds form a state machine:

\[
  \textsf{Setup} \to (\textsf{Registration})^* \to \textsf{Ring}
    \to (\textsf{Ballot})^* \to \textsf{Closed} \to \textsf{Tally}?
\]

\textsf{Setup} pins the election parameters: a 32-byte random
\texttt{election\_id} that doubles as the OTRS \emph{issue} value, the
manager's Ed25519 public key, ballot options, and the schedule
(\texttt{registration\_close} $\le$ \texttt{voting\_open} $\le$
\texttt{voting\_close}). \textsf{Registration} contains a voter's
Ristretto255 OTRS public key, a human-readable handle, and an Ed25519
\emph{attestation} by the manager over
$\texttt{voter-attestation-v1}\,\|\,\texttt{election\_id}\,\|\,\mathit{pk}\,\|\,\texttt{handle}$.
\textsf{Ring} is the canonical-order list of attested voter pks; after it
is published, registration is closed and voting opens. \textsf{Ballot}
carries an OTRS signature over $\texttt{ballot-v1}\,\|\,\texttt{choice}$,
bound to the issue and the ring. \textsf{Closed} is a marker.
\textsf{Tally} is the manager's claimed result --- which any auditor can
independently recompute.

\subsection{Witness federation}
\label{sec:witness}

The cohort signs entries; \emph{witnesses} sign \emph{heads}.
A witness is an independent third party with their own Ed25519 keypair.
Periodically (in the artifact, on-demand via the CLI) each witness
fetches the bulletin board, verifies it under the cohort, and emits a
\emph{checkpoint}
\[ \mathsf{Checkpoint} =
   (\mathit{log\_index},\ \mathit{head\_hash},\ \mathit{witness\_index},\ \sigma_{\mathsf{Ed}}) \]
signed under the witness's key with the domain-separated preimage
\verb+otrs-witness-checkpoint-v1 || log_index || head_hash || witness_index+.
Checkpoints accumulate in a sidecar file \texttt{<log>.witnesses} that
travels alongside the log.

Witnesses give us equivocation resistance: if a corrupted cohort
($\geq t$ malicious members) publishes two divergent logs, honest
witnesses who see both will sign two different \texttt{head\_hash}
values for the same \texttt{log\_index}. The auditor surfaces this as
hard \emph{equivocation evidence} — a publicly verifiable proof that the
cohort split. The genesis \textsf{ElectionSetup} pins the witness set
and a $k$-of-$M$ threshold; the auditor requires at least one log index
to be co-signed by $\geq k$ distinct witnesses (all agreeing on the
hash) before declaring the election trustworthy.

\subsection{Principals}

\textbf{Publisher cohort} (\texttt{voting/manager.py}). The $N$ cohort
members jointly publish \textsf{Setup}, register voters, freeze the
\textsf{Ring}, accept ballots, close voting, and publish a
\textsf{Tally}. Any subset of size $\geq t$ may sign each entry, either
synchronously (single-process tests) or asynchronously through the
pending-entry sidecar (\texttt{<log>.pending}) for the multi-process /
multi-machine deployment. The cohort is \emph{not} trusted to compute
the tally honestly: any tally they publish is a claim that auditors
recompute and refute.

\textbf{Voter} (\texttt{voting/voter.py}). Generates a Ristretto255 OTRS
keypair, sends the public key to the cohort out-of-band for inclusion,
fetches the bulletin board, and, once the ring is published, signs a
ballot for their chosen option using OTRS \Sign{} and submits it.

\textbf{Witnesses} (\texttt{voting/witness.py}). Each holds an Ed25519
key, fetches and verifies the log, and emits a checkpoint co-signing
the current head. Witnesses are \emph{read-only} — they cannot append
to the main log — and their security guarantee is that as long as $k$
of them are honest, equivocation by the cohort is publicly detectable.

\textbf{Auditor} (\texttt{voting/auditor.py}). Anyone with the bulletin
board can run the auditor; the cohort and witness identities are
bootstrapped from the genesis \textsf{Setup} entry itself, so the only
out-of-band trust is the channel by which an auditor obtained that
entry (typically a printed-on-a-poster cohort identity or a well-known
URL). The auditor:
\begin{enumerate}[leftmargin=*]
  \item verifies the chain (indices contiguous, prev-hashes match, $\geq t$
    valid cohort signatures per entry, timestamps non-decreasing);
  \item enforces the state machine, rejecting out-of-order records;
  \item requires the published ring to equal the multiset of voter pks
    appearing in cohort-attested \textsf{Registration} entries;
  \item OTRS-verifies every ballot against the issue, message, and ring;
  \item if witnesses are declared, verifies their checkpoints
    (signatures valid, no equivocation, $\geq k$ distinct witnesses
    co-sign at least one log index);
  \item groups ballots by signer using the OTRS trace columns
    (\cref{sec:construction}), counts a single vote per voter, and
    excludes ballots from any voter who signed two distinct messages;
  \item compares the canonical tally against the cohort's \textsf{Tally}
    record, if present.
\end{enumerate}
The CLI exposes the auditor as \texttt{evote audit} and exits non-zero
on any failure or tally mismatch.

\section{Architecture B: proof-of-work chain}
\label{sec:chain}

\texttt{chain/} implements the same election lifecycle on top of a
Bitcoin-shaped proof-of-work chain. The cryptographic primitive
(\cref{sec:construction}) and the record schema
(\texttt{voting/records.py}, \cref{sec:records}) are reused
\emph{verbatim}. What changes is the storage and ordering layer: instead
of a $t$-of-$N$ cohort co-signing every entry, anyone may mine, and the
canonical history is the chain with the greatest cumulative
$2^{\text{difficulty}}$ (GHOST-style heaviest-chain selection).

\subsection{Block and chain format}

A block has a 121-byte fixed header
\[ (\mathit{prev\_hash},\ \mathit{height},\ \mathit{difficulty},\
   \mathit{timestamp},\ \mathit{miner\_pk},\ \mathit{tx\_root},\
   \mathit{nonce}) \]
plus a variable-length transaction list. The PoW puzzle is the standard
SHA-256-of-header inequality:
\[ \mathsf{leading\_zero\_bits}(\mathrm{SHA{-}256}(\mathit{header}))
   \geq \mathit{difficulty}. \]
Transactions are bound into the header via $\mathit{tx\_root} =
\mathrm{SHA{-}256}(\Vert_{k} \mathrm{canon}(\mathit{tx}_k))$ (a flat
hash rather than a Merkle root --- block bodies are small enough that
inclusion proofs are unnecessary in the prototype). Difficulty is
adjusted $\pm 1$ when the median inter-block interval over a sliding
window of 11 blocks falls outside $[0.5T,\ 1.5T]$ around a target
$T = 60$~seconds, mirroring the legacy node's adjustment rule.

\subsection{Account ledger and transactions}
\label{sec:chain-tx}

Every Ed25519 public key maps to an account with a \emph{balance} and a
\emph{nonce}. Eight transaction kinds are defined; six wrap the
\texttt{voting/records.py} record types so the per-record validation
already proved correct in \cref{sec:records} is reused:

\begin{center}\small
\begin{tabular}{lll}
\toprule
Kind & Signed by & Wraps \\
\midrule
\texttt{coinbase}        & none (per-block, first tx) & --- (mints reward + fees) \\
\texttt{transfer}        & sender                     & --- (eVote payment) \\
\texttt{setup\_poll}     & poll creator               & \textsf{ElectionSetup} \\
\texttt{register\_voter} & sponsor                    & \textsf{VoterRegistration} \\
\texttt{publish\_ring}   & poll creator               & \textsf{RingPublication} \\
\texttt{ballot}          & \emph{none}                & \textsf{Ballot} \\
\texttt{close\_poll}     & poll creator               & \textsf{VotingClosed} \\
\texttt{tally}           & poll creator               & \textsf{TallyPublication} \\
\bottomrule
\end{tabular}
\end{center}

The \texttt{ballot} kind is the only transaction with \emph{no} chain-layer
sender, signature, or fee. This is essential for anonymity: a fee-bearing
ballot would expose the voter's Ed25519 identity on chain. Validity is
the OTRS signature inside the payload, verifiable by anyone against the
published ring. Outside the ring, no valid signature can be produced, so
flooding the mempool with bogus ballots is computationally bounded by
ring membership and rejected at validation; miners are expected to
include all valid ballots they see.

\subsection{State machine}

The same lifecycle as Architecture~A is enforced at the chain layer.
\texttt{chain.state} maintains a per-poll \texttt{PollState} indexed by
\texttt{election\_id}, with phase transitions
$\textsf{registration} \to \textsf{voting} \to \textsf{closed} \to
\textsf{tallied}$ triggered by the corresponding transaction kinds.
Per-tx checks include: poll-creator identity match
(only the address that issued the \texttt{setup\_poll} may issue ring /
close / tally), per-poll registration uniqueness, ring membership
equals the multiset of attested registrations, ballot timestamp
within $[\textsf{voting\_open},\ \textsf{voting\_close}]$, and OTRS
verification of every ballot at the moment it is included in a block.

Block application is atomic: a snapshot of the state is taken at entry,
mutations happen on the snapshot, and the snapshot is committed only on
success --- so a malformed transaction does not leave the chain in a
partial state.

\subsection{Mining and node lifecycle}

The miner (\texttt{chain/mining.py}) takes a header template (with
\texttt{nonce} = 0) and scans nonces until the PoW threshold is met.
For the artifact prototype this is single-process; a production miner
would partition the nonce space across cores or machines.

\texttt{chain.node.Node} owns the canonical chain and the folded state;
\texttt{append(block)} validates and applies; \texttt{save/load}
persist the chain as a length-prefixed binary stream and rebuild state
from scratch on load (so a corrupted state cache cannot silently
desync from chain content). Fork resolution by cumulative weight is
\emph{not} implemented in this prototype --- the node accepts the chain
it is given; the necessary information (\texttt{State.cum\_weight}) is
available and a heaviest-chain switch is a small extension.

\subsection{Audit}
\label{sec:chain-audit}

\texttt{chain.auditor.audit\_chain} loads the chain through
\texttt{Node.load}, which already enforces PoW, signatures, nonces,
balances, the state machine, and OTRS-per-ballot validity. The auditor
then runs the same $\sigma$-column bucketing as
\texttt{voting.auditor.audit} (\cref{sec:bb}) to compute the canonical
tally, identify double-signers, and compare against any on-chain
\textsf{TallyPublication}. The tally algorithm is byte-identical
between the two architectures because it depends only on the ballots
and the ring, not on how either of them got there.

\subsection{CLI walk-through}

Each subcommand that produces a new state mines its transaction into a
fresh block on top of the current tip --- one CLI call equals one
block. This is simpler than a mempool for demo purposes and makes the
demonstration auditable line-by-line.

\begin{lstlisting}[style=shell]
python3 -m chain.cli account-keygen --out alice.json
python3 -m chain.cli voter-keygen   --out v0.json
python3 -m chain.cli init-chain --chain c.bin --miner miner.json \
    --allocate "$(jq -r .pk_b64 alice.json):1000" --timestamp $T0
EID=$(python3 -m chain.cli setup-poll --chain c.bin --miner miner.json \
    --creator alice.json --title "Demo" --options "yes,no" \
    --registration-close $T1 --voting-open $T2 --voting-close $T3)
python3 -m chain.cli register-voter --chain c.bin --miner miner.json \
    --sponsor alice.json --election-id $EID \
    --voter-pk $(jq -r .pk_b64 v0.json) --handle v0
python3 -m chain.cli publish-ring --chain c.bin --miner miner.json \
    --creator alice.json --election-id $EID
python3 -m chain.cli vote --chain c.bin --miner miner.json \
    --voter v0.json --election-id $EID --choice yes
python3 -m chain.cli close-poll --chain c.bin --miner miner.json \
    --creator alice.json --election-id $EID
python3 -m chain.cli audit --chain c.bin
\end{lstlisting}

The audit prints the chain head, the per-poll tally, and any
double-sign culprits, and exits non-zero on any inconsistency.

\section{Architecture comparison}
\label{sec:compare}

The two architectures share an entire cryptographic stack --- the same
OTRS primitive, the same record schema, the same canonical encodings,
the same audit-side algorithm for tally and trace. What differs is the
storage and ordering layer, and that single delta cascades into the
operational properties below.

\begin{table}[h]
\centering
\caption{Architecture A (federated bulletin board) vs.\ B (PoW chain).}
\label{tab:compare}
\small
\begin{tabularx}{\textwidth}{l X X}
\toprule
Property & \textbf{A: federated} (\texttt{voting/}) & \textbf{B: PoW chain} (\texttt{chain/}) \\
\midrule
Validator membership   & permissioned: $N$ named publishers at genesis & permissionless: anyone with a miner \\
Liveness               & against $N{-}t$ unavailable publishers       & against any minority hashpower \\
Safety from forks      & threshold-signed total order; no forks       & probabilistic, heaviest chain wins \\
Censorship resistance  & any $t$ honest publishers can publish        & any miner can include any valid tx \\
Equivocation defence   & witness federation, $k$-of-$M$ checkpoints   & longest-chain rule; reorgs always possible \\
Sybil resistance for voters & cohort-signed registration & poll-creator-signed registration (same) \\
Cost per ballot        & 1 SHA-256 + threshold Ed25519 sigs           & 1 block of PoW + tx validation \\
Wall-clock to close    & seconds                                      & minutes (target block interval) \\
Disenfranchisement risk& zero (voters do not transact)                & non-zero (sponsor or self pays fees) \\
Cross-poll linkability & per-election ring; no global identity        & persistent Ed25519 addresses for poll creators / sponsors; voter OTRS keys can be rotated per poll \\
Receipt-freeness       & open (same in both)                          & open (same in both) \\
Required infrastructure& $N + M$ named parties + their key custody    & miners + an eVote bootstrap distribution \\
Failure mode if abandoned & log goes stale, audit still works         & chain stops; abandoned blocks rewritable by anyone with hashpower \\
Lines of code          & $\approx$1\,800 + 111 tests                  & $\approx$1\,400 + 16 tests \\
\bottomrule
\end{tabularx}
\end{table}

\subsection{When each architecture wins}

\textbf{Architecture A wins when} the election has a known, accountable
authority set (a university senate, a corporate board, an
inter-institution treaty body), when ballots need to be finalised in
seconds rather than minutes, and when the operational footprint must
not exceed the named cohort plus a handful of independent witnesses.
It is essentially the Certificate-Transparency / Sigsum / sigstore
pattern applied to elections: small named log + threshold trust +
independent witnesses for equivocation defence. No tokens, no fees,
no mining, no chain reorganisation.

\textbf{Architecture B wins when} permissionless validator membership
is a non-negotiable design constraint (e.g.\ no party --- including
the poll creator --- is trusted to host the canonical record), when
many polls are expected to share the same infrastructure and pay for
its operation, and when the operator is comfortable accepting
probabilistic finality and a non-trivial cost-per-vote. The chain has
no genesis-pinned trust set beyond \emph{the rules of the protocol}.

\subsection{When neither wins}

Both architectures expose persistent identities at the chain layer:
in A the cohort and witness public keys, in B the poll creator and
sponsor public keys (and any Ed25519 addresses used to fund eVotes).
Neither architecture solves \emph{coercion resistance} --- a voter
who reveals their OTRS secret to a coercer enables trace-tag matching
in either system. Neither defends against \emph{compromised voter
devices} --- a malicious client trivially signs ballots the user did
not authorise. Both gaps are independent of the storage layer; they
sit upstream of any OTRS-based voting design and require
JCJ/Civitas-style coercion-resistance machinery to close.

\subsection{The economic asymmetry, addressed honestly}

In A the security budget is the operational cost of running $N$
publishers plus $M$ witnesses --- a one-off setup expense, bounded.
In B the security budget per block is the fee-plus-reward paid to
miners, which over a window is bounded by aggregate poll fees on the
chain. Crucially, the \emph{attacker's} payoff scales with the largest
poll on the chain at any moment, not with average fees. For elections
of meaningful political stakes (corporate boards, public referenda) a
sufficiently funded adversary can rent hashpower for the voting window
at a cost below their expected gain from controlling
outcome~\cite{ParkSpecter2021}. Architecture~A sidesteps this entirely
because it does not try to be Sybil-resistant via economic incentive
--- it relies on \emph{naming} the trust set and making cheating by
that set publicly evident. Whether that trade is acceptable is a
deployment-specific judgment, not a cryptographic one.

\section{Threat model (summary)}
\label{sec:threat}

The full threat model lives in \texttt{paper/threat\_model.md}. Briefly,
under DDH in Ristretto255, the random-oracle model for SHA-512--based
hash-to-curve, and Ed25519 unforgeability, v0.3 provides:
\textit{(i)} anonymity inside the ring for any single ballot;
\textit{(ii)} public, deterministic detection of double-voting within an
election; \textit{(iii)} publicly recomputable tallying;
\textit{(iv)} tamper-evident history; \textit{(v)} eligibility
enforcement against the cohort's attestation set;
\textit{(vi)} liveness against $N - t$ unavailable cohort members;
\textit{(vii)} soundness against $t - 1$ corrupted cohort members; and,
when witnesses are configured, \textit{(viii)} publicly verifiable
equivocation evidence whenever the cohort itself splits. It does
\emph{not} provide receipt-freeness (a voter who reveals $x_i$ to a
coercer enables trace-tag matching), defence against compromised voter
devices, or post-quantum security. These remaining gaps are the v0.4
roadmap.

\section{Cryptographic security}
\label{sec:security}

The scheme is proved unforgeable, anonymous, and traceable in the random
oracle model under DDH~\cite{ScafuroZhang2021}. We do not reproduce the
proofs here; we restate the security claims as the implementation enforces
them and flag where our instantiation deviates from the paper's abstract
model.

\begin{itemize}[leftmargin=*]
\item \textbf{Unforgeability} follows from soundness of the $\Sigma$-OR
      proof: producing a valid $(\{c_j\}, \{z_j\})$ tuple with
      $\sum_j c_j \equiv \Htwo(\cdot)$ without knowing any $x_j$ requires
      inverting $\Htwo$ or solving discrete logarithm.
\item \textbf{Anonymity} follows from honest-verifier zero-knowledge of
      each individual $\Sigma$-proof, plus DDH for the trace tag $\sigma_i$.
\item \textbf{Traceability} follows because $\sigma_{i^*} = h^{x_{i^*}}$ is
      fully determined by $(I, R, x_{i^*})$ and independent of $m$.
\end{itemize}

Our instantiation matches the abstract assumptions modulo two choices: the
hash-to-curve is \emph{modelled} as a random oracle (XMD:SHA-512 RO is the
cryptographic-community-standard instantiation), and the group is
prime-order Ristretto255 rather than an abstract prime-order group (a
strictly safer instantiation than the original DSA subgroup, ruling out
small-subgroup and cofactor attacks by construction).

\subsection{Known limitations}
\label{sec:limitations}

\begin{itemize}[leftmargin=*]
\item \textbf{Constant-time.} Python integer arithmetic is not
      constant-time; the point--scalar product delegates to libsodium which
      \emph{is} constant-time. The remaining timing channel is the
      \texttt{is\_zero} branch in \texttt{Point.scalar\_mul}, which is
      reached only with probability $\approx 2^{-252}$ for random scalars
      derived in our protocol---negligible.
\item \textbf{Equality of group elements} uses Python \texttt{bytes}
      equality, which is not constant-time. Group-element equality is
      \emph{not} secret-dependent in our verification path (both sides are
      derived from public inputs), so this is acceptable.
\item \textbf{No HSM-style binding} between a signer's secret and the host:
      a compromised client trivially signs anything.
\end{itemize}

\begin{remark}[$n = 1$ is degenerate]
The trace algorithm's ``all columns match'' and ``exactly one column matches''
predicates coincide for a ring of size~1. In that case \Trace{} returns
``linked'' even on a double-sign attempt. Anonymity is undefined for $n = 1$
regardless, so callers should require $n \ge 2$ for any realistic use.
\end{remark}

\section{Evaluation}
\label{sec:eval}

Microbenchmarks were collected on a single core (no parallelism) on a Linux
laptop, with five repetitions per operation, median reported.

\subsection{Cryptographic primitive}

Reproduce with:
\begin{lstlisting}[style=shell]
python3 -m bench.bench_otrs --sizes 2,4,8,16,32,64,128 --output bench/results.csv
\end{lstlisting}

\begin{table}[h]
\centering
\caption{OTRS cost vs.\ ring size $n$.}
\label{tab:bench}
\begin{tabular}{rrrrr}
\toprule
$n$ & Sign (ms) & Verify (ms) & Trace (ms) & Sig size (B) \\
\midrule
2   & 0.67  & 0.72  & 0.42  & 164  \\
4   & 1.77  & 1.72  & 0.68  & 292  \\
8   & 2.62  & 3.14  & 1.32  & 548  \\
16  & 5.20  & 5.79  & 2.57  & 1\,060 \\
32  & 10.42 & 10.51 & 5.27  & 2\,084 \\
64  & 20.77 & 20.89 & 10.77 & 4\,132 \\
128 & 45.02 & 42.70 & 20.37 & 8\,228 \\
\bottomrule
\end{tabular}
\end{table}

Cost is linear in $n$---empirically about $0.34$~ms of sign time and
$0.33$~ms of verify time per ring member at this configuration, with the
constant dominated by the $2n + 1$ scalar multiplications each side
performs. Trace runs in $\approx 0.16$~ms per member because it computes
only $n$ scalar multiplications per signature rather than $2n + 1$
exponentiations plus group additions. Signature size is
$|A_1| + 4 + 2n \cdot |\mathrm{scalar}| = 32 + 4 + 64n$~bytes, matching
the measured column to the byte.

\subsection{Full election}

The end-to-end cost of an election is dominated by the OTRS work: each
ballot needs one \Sign{} call (signer side, $O(n)$) and one \Verify{}
call (auditor side, $O(n)$). The bulletin-board layer adds only a SHA-256
hash, an Ed25519 signature, and a file append per record, all in the
microsecond range. \cref{tab:syseval} shows the wall-clock cost of a
complete election for varying voter counts $N$, where every registered
voter casts one ballot and the auditor verifies the entire log.

\begin{table}[h]
\centering
\caption{Full election cost (single core, ring size = number of voters).}
\label{tab:syseval}
\begin{tabular}{rrrrrr}
\toprule
$N$ & Register (ms) & Ring publish (ms) & All ballots (ms) & Audit (ms) & Log size (KB) \\
\midrule
10  & 6.4   & 0.9 & 41.3    & 51.2    & 21.2  \\
25  & 15.1  & 0.9 & 245.6   & 277.7   & 93.0  \\
50  & 45.1  & 1.5 & 1\,004.7 & 1\,104.0 & 323.8 \\
\bottomrule
\end{tabular}
\end{table}

Ballot submission and audit are both $O(N^2)$ in this configuration
because the ring grows with $N$ and the per-ballot cost is $O(\text{ring
size})$. For larger elections one would either (i) shard voters into
multiple parallel rings, or (ii) move to a logarithmic-size traceable
ring (\cref{sec:future}, item~1). At $N \le 50$ the entire end-to-end
flow finishes in under $1.2$ seconds on a single core; the bulletin-board
log fits in a third of a megabyte.

\section{Comparison to related work}
\label{sec:related}

A spectrum of ring-signature constructions exists, parametrised by
signature size, signer/verifier cost, trust assumptions, and traceability
properties.

\begin{table}[h]
\centering
\caption{Ring-signature design space.}
\label{tab:related}
\small
\begin{tabularx}{\textwidth}{l l l l X}
\toprule
Scheme & Sig size & Trace & Setup & Assumption \\
\midrule
RST01~\cite{RST01} & $O(n)$ & none & none & RSA-like \\
LSAG/CLSAG~\cite{CLSAG2020} & $O(n)$ & linkable & none & DDH \\
MLSAG~\cite{Noether2015} & $O(n)$ & linkable & none & DDH \\
Triptych~\cite{Triptych2020} & $O(\log n)$ & linkable & none & DDH \\
Lelantus~\cite{Jivanyan2019} & $O(\log n)$ & linkable & none & DDH \\
Raptor / Falafl~\cite{Raptor2018} & $O(\log n)$ & linkable & none & Lattice (PQ) \\
\textbf{Scafuro--Zhang}~\cite{ScafuroZhang2021} & $O(n)$ & \textbf{traceable (one-time)} & none & DDH, ROM \\
\bottomrule
\end{tabularx}
\end{table}

The relevant axis for our use case (one ballot per voter, public
auditability, no trusted authority) is \emph{traceability}, which is rarer
than linkability. Most linkable schemes only reveal that two signatures
share a signer without identifying them; OTRS additionally identifies the
misbehaving public key. The cost is signature size: $O(n)$ versus
$O(\log n)$ for state-of-the-art linkable schemes.

A natural research direction (\cref{sec:future}) is to combine the trace tag
with a log-size set-membership proof \`a la Groth--Kohlweiss~\cite{GK15} or
Triptych.

\section{Open problems and future work}
\label{sec:future}

\begin{enumerate}[leftmargin=*]
\item \textbf{Logarithmic-size traceable rings.} Replace the linear
      $\Sigma$-OR with a Groth--Kohlweiss-style membership
      proof~\cite{GK15} while keeping the trace tag $\sigma_i = h^{x_i}$.
      The challenge is binding the per-position $A_1$ reconstruction to the
      index that the membership proof commits to.
\item \textbf{Formal verification.}
      \begin{itemize}
        \item \emph{EasyCrypt}: model the three games (unforgeability,
              anonymity, traceability) and machine-check the reductions to
              DDH and ROM. The scheme's small algebraic surface makes this
              tractable.
        \item \emph{ProVerif/Tamarin}: model the \emph{protocol layer} of
              the e-voting application---registration, ballot submission,
              tracing---and verify ballot anonymity and double-vote
              detection as observational equivalence properties.
      \end{itemize}
\item \textbf{Post-quantum variants.} The DDH reduction breaks under Shor's
      algorithm. Lattice-based linkable rings (Raptor, Calamari, Falafl)
      exist; constructing a \emph{traceable} lattice ring with one-time
      tags is, to our knowledge, open. A literature survey plus barrier
      analysis is a tractable Master's thesis topic.
\item \textbf{Coercion resistance.} OTRS prevents \emph{vote duplication}
      but not \emph{vote selling}---a coerced voter could prove to a
      coercer which signature is theirs by revealing $x_i$. A receipt-free
      extension \`a la JCJ~\cite{JCJ2005} or Civitas~\cite{Civitas2008} is
      needed for coercion resistance. This is the single most important
      gap in v0.2; see \texttt{paper/threat\_model.md} \S 4.
\item \textbf{True threshold Ed25519 (FROST).} v0.3's $t$-of-$N$
      threshold is implemented as a vector of $t$ separate Ed25519
      signatures, one per signer. This is the convention used by CT and
      Sigsum and is sound, but a single FROST~\cite{FROST2020}
      aggregated signature would compress entries from $64t$ to $64$
      bytes. A clean follow-up is to swap the vector for FROST without
      changing the rest of the system.
\item \textbf{Gossip protocol for bulletin board distribution.} v0.3
      treats the log and the witness sidecar as local files. A
      production deployment would replicate them across the cohort and
      a wider set of mirrors, with auditors fetching from $\geq 2$
      independent mirrors and cross-checking. This is engineering work
      rather than research, and is independent of the cryptography.
\item \textbf{Cross-validation against RFC~9380 test vectors} for
      ristretto255 hash-to-curve; we ship regression vectors but not (yet)
      the official ones, since RFC~9380's appendix omits ristretto255
      vectors in some normative drafts.
\end{enumerate}

\section{Reproducibility}
\label{sec:repro}

This artifact is released under the MIT license at
\url{https://github.com/NickP005/e-ring-voting}.

\begin{lstlisting}[style=shell]
# system deps (Debian/Ubuntu)
sudo apt install -y libsodium23 python3-cffi python3-pytest \
    python3-hypothesis python3-cryptography

# full test suite: 127 tests across both architectures + crypto + properties
python3 -m pytest -v

# OTRS benchmarks
python3 -m bench.bench_otrs --sizes 2,4,8,16,32,64,128

# Architecture A: federated bulletin-board demo
python3 -m voting.cli manager-keygen --sk-out msk.pem --pk-out mpk.pem
python3 -m voting.cli setup --log log.jsonl --sk msk.pem --title "Demo" \
    --options "yes,no,abstain" \
    --registration-close <T0> --voting-open <T1> --voting-close <T2>
# ...voter-keygen / register / publish-ring / vote / close / tally...
python3 -m voting.cli audit --log log.jsonl --pk mpk.pem

# Architecture B: PoW chain demo
python3 -m chain.cli account-keygen --out alice.json
python3 -m chain.cli voter-keygen   --out v0.json
python3 -m chain.cli init-chain --chain c.bin --miner alice.json \
    --allocate "$(jq -r .pk_b64 alice.json):1000" --timestamp <T0>
# ...setup-poll / register-voter / publish-ring / vote / close-poll...
python3 -m chain.cli audit --chain c.bin
\end{lstlisting}

\begin{thebibliography}{99}\small
\bibitem{Helios2008} B.~Adida. Helios: Web-based Open-Audit Voting. USENIX Security, 2008.
\bibitem{ElectionGuard2021} J.~Benaloh et al. ElectionGuard: a Cryptographic Toolkit to Enable Verifiable Elections. 2021.
\bibitem{Scantegrity2008} D.~Chaum et al. Scantegrity II. EVT, 2008.
\bibitem{Civitas2008} M.~Clarkson, S.~Chong, A.~Myers. Civitas: Toward a Secure Voting System. IEEE S\&P, 2008.
\bibitem{Ristretto2020} H.~de Valence et al. The Ristretto Group. \url{https://ristretto.group}.
\bibitem{FS07} E.~Fujisaki, K.~Suzuki. Traceable Ring Signature. PKC, 2007.
\bibitem{GK15} J.~Groth, M.~Kohlweiss. One-Out-of-Many Proofs. EUROCRYPT, 2015.
\bibitem{CLSAG2020} B.~Goodell et al. Concise Linkable Ring Signatures and Forgery Against Adversarial Keys. 2020.
\bibitem{Jivanyan2019} A.~Jivanyan. Lelantus: Toward Confidentiality and Anonymity of Blockchain Transactions. 2019.
\bibitem{JCJ2005} A.~Juels, D.~Catalano, M.~Jakobsson. Coercion-resistant electronic elections. WPES, 2005.
\bibitem{LWW04} J.~K.~Liu, V.~K.~Wei, D.~S.~Wong. Linkable Spontaneous Anonymous Group Signature for Ad Hoc Groups. ACISP, 2004.
\bibitem{Noether2015} S.~Noether. Ring Signature Confidential Transactions for Monero. 2015.
\bibitem{Triptych2020} S.~Noether et al. Triptych: Logarithmic-Sized Linkable Ring Signatures. 2020.
\bibitem{Raptor2018} X.~Lu, M.~H.~Au, Z.~Zhang. Raptor: A Practical Lattice-Based (Linkable) Ring Signature. ACNS, 2019.
\bibitem{RFC9380} S.~Scott, S.~Wood, R.~Wahby, A.~Faz-Hern\'andez, N.~Sullivan. Hashing to Elliptic Curves. RFC~9380, 2023.
\bibitem{RST01} R.~Rivest, A.~Shamir, Y.~Tauman. How to Leak a Secret. ASIACRYPT, 2001.
\bibitem{Springall2014} D.~Springall et al. Security Analysis of the Estonian Internet Voting System. CCS, 2014.
\bibitem{ScafuroZhang2021} A.~Scafuro, B.~Zhang. One-time Traceable Ring Signatures. ESORICS, 2021.
\bibitem{FROST2020} C.~Komlo, I.~Goldberg. FROST: Flexible Round-Optimized Schnorr Threshold Signatures. SAC, 2020.
\bibitem{ParkSpecter2021} S.~Park, M.~Specter et al. Going from bad to worse: from Internet voting to blockchain voting. Journal of Cybersecurity, 2021.
\end{thebibliography}

\end{document}
